\documentclass[10pt,a4paper]{article}
\usepackage[margin=1.1cm]{geometry}
\usepackage{amsmath,amsthm,amsfonts,amssymb,amscd,cite,graphicx}

\usepackage{titlesec}
\titleformat{\section}
{\normalfont\fontsize{12}{15}\bfseries}{\thesection}{1em.}{}

\usepackage[labelfont=bf]{caption}
\counterwithin{figure}{section}
\counterwithin{table}{section}

\newtheorem{proposition}{Proposition}[section]
\newtheorem{conjecture}{Conjecture}[section]
\newtheorem{corollary}{Corollary}[section]
\newtheorem{lemma}{Lemma}[section]
\newtheorem{definition}{Definition}[section]
\newtheorem{theorem}{Theorem}[section]
\newtheorem{remark}{Remark}[section]
\newtheorem{example}{Example}[section]
\newtheorem{counterexample}{Counter Example}[section]
\newtheorem{observation}{Observation}[section]

\renewcommand{\thefootnote}{\fnsymbol{footnote}}
\allowdisplaybreaks[4]

\let\oldbibliography\thebibliography
\renewcommand{\thebibliography}[1]{%
  \oldbibliography{#1}%
  \setlength{\itemsep}{-2pt}%
}

\baselineskip=1.20in

\begin{document}

\baselineskip=0.20in

\makebox[\textwidth]{%
\hglue-15pt
\begin{minipage}{0.6cm}
\vskip9pt
\end{minipage} \vspace{-\parskip}
\begin{minipage}[t]{6cm}
\footnotesize{ {\bf Discrete Mathematics Letters} \\ \underline{www.dmlett.com}}
\end{minipage}
\hfill
\begin{minipage}[t]{6.5cm}
\normalsize {\it Discrete Math. Lett.}  {\bf X} (202X) XX--XX
\end{minipage}}
\vskip36pt

\noindent
{\large \bf A Parity Lower Bound for Twofold Radius-One Coverings of the Binary Hypercube}\\

\noindent
Hanyu Yang\footnote{E-mail address: hanyu.yang.92@gmail.com}\\

\noindent
\footnotesize {\it Fremont, California, United States\/}\\

\noindent
 (\footnotesize Received: Day Month 202X. Received in revised form: Day Month 202X. Accepted: Day Month 202X. Published online: Day Month 202X.)\\

\setcounter{page}{1} \thispagestyle{empty}

\baselineskip=0.20in

\normalsize

\begin{abstract}
\noindent
A binary code is a twofold covering of radius one if every word of the
hypercube lies within Hamming distance one of at least two codewords. For every
even length we prove a lower bound on the minimum size of such a code by a
parity observation taken at a codeword rather than at an arbitrary word: in a
code without repeated words the ball of a codeword carries a unique diagonal
term, whose parity forces an excess whenever the length is even. After taking
ceilings, the bound exceeds the general lower bound of Krotov and Potapov
for every even length from six onwards. It raises the tabulated lower bounds at the five
lengths eight, ten, twelve, fourteen and sixteen, and at length six it supplies
an elementary lower-bound proof which, with the known twenty-word construction,
recovers the exact value. The argument extends to every feasible even
multiplicity; for the underlying real-valued bounds we determine exactly when
the resulting expression exceeds the Krotov--Potapov expression, noting that a
strict gain need not survive the ceiling.
\\[2mm]
{\bf Keywords:} covering code; multiple covering; twofold covering; double domination; covering excess; hypercube; lower bound.\\[2mm]
{\bf 2020 Mathematics Subject Classification:} 94B65, 05B40, 94B75.
\end{abstract}

\baselineskip=0.20in

\section{Introduction}

Let $\mathbb{F}_2^n$ be the binary Hamming space and let
$B(x)=\{y: d(x,y)\le 1\}$ be the closed ball of radius one, so $|B(x)|=n+1$.
A set $C\subseteq\mathbb{F}_2^n$ is a \emph{$\mu$-fold covering} of radius one
if $|B(x)\cap C|\ge\mu$ for every $x\in\mathbb{F}_2^n$, and $K(n,1,\mu)$ denotes
the least cardinality of such a set. Codewords are distinct: $C$ is a set and
not a multiset. Allowing repetitions gives a different quantity, treated
separately in \cite{HHKL-1993}, and the hypothesis that $C$ is a set is used
essentially below. We write $K(n,1,2)$ for the twofold case. Equivalently, such
a set is a $\mu$-tuple dominating set in the $n$-dimensional hypercube graph,
and a double dominating set when $\mu=2$.

The standard tool for lower bounds of this kind is the covering excess method,
extended to multiple coverings by H\"am\"al\"ainen, Honkala, Kaikkonen and
Litsyn \cite{HHKL-1993}. Their Theorem 6 improves on the sphere covering bound
by an additive correction term $\varepsilon$ which, at radius one, vanishes
exactly when $\mu(n+1)$ is even. For even $n$ that is exactly when $\mu$ is
even, so at $\mu=2$ and $n$ even their bound reduces to the sphere covering
bound; see Section \ref{sec-prior}.

Our observation is that a \emph{different} parity is available, and that it
fires precisely where theirs vanishes. Summing the coverage function over the
ball of a \emph{codeword} rather than of an arbitrary word produces a sum with
a unique diagonal term $n+1$, hence a sum that is odd whenever $n$ is even,
forcing an excess in every codeword ball. The resulting bound is Theorem
\ref{thm-main}.

\begin{theorem}\label{thm-main}
Let $n\ge 2$ be even and let $C\subseteq\mathbb{F}_2^n$ satisfy
$|B(x)\cap C|\ge 2$ for every $x\in\mathbb{F}_2^n$. Then
\[
  |C| \;\ge\; \frac{3\cdot 2^{\,n+1}}{3n+2},
  \qquad\text{and hence}\qquad
  K(n,1,2) \;\ge\; \left\lceil \frac{3\cdot 2^{\,n+1}}{3n+2} \right\rceil .
\]
\end{theorem}

The proof occupies half a page and uses nothing beyond double counting. Section
\ref{sec-dominance} shows that after taking ceilings the bound exceeds the
general lower bound of Krotov and Potapov \cite{KP-2021} at every even
$n\ge 6$; Section \ref{sec-small} records the consequences at small $n$; and
Section \ref{sec-mu} carries the argument to every feasible even multiplicity
and classifies when the underlying real-valued bound exceeds the
Krotov--Potapov bound.

\section{Notation}\label{sec-notation}

For $C\subseteq\mathbb{F}_2^n$ put
\[
  c(y)=|B(y)\cap C|,\qquad g(y)=c(y)-2,\qquad E=\sum_{y\in\mathbb{F}_2^n} g(y).
\]
If $C$ is a twofold covering then $g\ge 0$ pointwise. Counting the incidences
$(y,z)$ with $z\in C$ and $y\in B(z)$ in two ways gives
$\sum_y c(y)=(n+1)|C|$, whence
\begin{equation}\label{eq-excess}
  E=(n+1)|C|-2^{\,n+1}.
\end{equation}
Let $S=\{y: g(y)\ge 1\}$ be the \emph{over-covered set} and put $s=|S|$. Since
$g$ is non-negative and integer-valued, and vanishes off $S$,
\begin{equation}\label{eq-s-le-E}
  s\le E.
\end{equation}
Finally we use the ball intersection numbers of the cube: for $x\ne z$,
\begin{equation}\label{eq-ball-int}
  |B(x)\cap B(z)| =
  \begin{cases}
    2, & d(x,z)\in\{1,2\},\\
    0, & d(x,z)\ge 3,
  \end{cases}
\end{equation}
together with $|B(x)\cap B(x)|=n+1$. For $d(x,z)=1$ the two common points are
$x$ and $z$ themselves; for $d(x,z)=2$ they are the two words lying between
them.

\section{Proof of Theorem \ref{thm-main}}\label{sec-proof}

\begin{proof}
\textbf{Step 1 (parity).} Fix $x\in C$ and sum the coverage function over its
ball. Exchanging the order of summation and applying \eqref{eq-ball-int},
\[
  \sum_{y\in B(x)} c(y)
  \;=\; \sum_{z\in C} |B(x)\cap B(z)|
  \;=\; (n+1) + 2N_2(x),
\]
where $N_2(x)=\#\{z\in C: 1\le d(x,z)\le 2\}$. The term $n+1$ is the
contribution of $z=x$, which lies in $C$ by assumption and, because $C$ is a
set, occurs exactly once; every other codeword contributes $0$ or $2$. Because
$n$ is even, $n+1$ is odd, so the whole sum is \emph{odd}. On the other hand it
is a sum of $n+1$ terms each at least $2$, so it is at least $2(n+1)$; being
odd, it is at least $2(n+1)+1$. Therefore
\[
  \sum_{y\in B(x)} g(y)\;\ge\;1 ,
\]
so every codeword $x$ admits some $y\in B(x)$ with $g(y)\ge 1$, that is, with
$y\in S$.

\textbf{Step 2 (incidence count).} Count the pairs $(x,y)$ with $x\in C$,
$y\in S$ and $y\in B(x)$. By Step 1 each $x\in C$ occurs at least once, and each
$y\in S$ occurs exactly $|B(y)\cap C|=c(y)$ times. Hence
\[
  |C| \;\le\; \sum_{y\in S} c(y) \;=\; \sum_{y\in S}\bigl(g(y)+2\bigr)
      \;=\; E + 2s .
\]

\textbf{Step 3.} By \eqref{eq-s-le-E} we have $s\le E$, so $|C|\le 3E$.
Substituting \eqref{eq-excess},
\[
  |C| \;\le\; 3(n+1)|C| - 3\cdot 2^{\,n+1},
  \qquad\text{that is}\qquad
  3\cdot 2^{\,n+1} \;\le\; (3n+2)\,|C| . \qedhere
\]
\end{proof}

\begin{remark}\label{rem-parity}
The parity step is essential, and the conclusion genuinely fails for odd $n$.
For example, in $\mathbb{F}_2^3$ the code $C=\{000,001,110,111\}$ has $c(x)=2$
for every $x$,
so it is a perfect twofold covering with $E=0$ and $|C|=4=K(3,1,2)$, whereas the
formula of Theorem \ref{thm-main} would demand $\lceil 48/11\rceil=5$.
\end{remark}

\begin{remark}\label{rem-set}
The hypothesis that $C$ is a set is essential for this parity argument. If a
codeword $x$ could occur with multiplicity $m_x$, the diagonal contribution in
Step 1 would be $m_x(n+1)$, whose parity is not determined when $m_x$ may be
even. The incidence count of Step 2 adapts to multiplicities; it is the
uniqueness of the diagonal term that fails.
\end{remark}

\section{Comparison with the Krotov--Potapov Bound}\label{sec-dominance}

Krotov and Potapov \cite[Theorem 6]{KP-2021} give, with $\mu\equiv\tau\pmod 2$
and $\tau\in\{0,1\}$,
\[
  K(n,1,\mu)\ge\frac{2^n(\mu n+3\mu+\tau)}{n(n+4)}\ \ (n\equiv 0\bmod 4),
  \qquad
  K(n,1,\mu)\ge\frac{2^n(\mu n+\mu+\tau)}{n(n+2)}\ \ (n\equiv 2\bmod 4).
\]
Write $L(n)=6\cdot 2^n/(3n+2)$ for the bound of Theorem \ref{thm-main}, and
$KP(n)$ for the relevant right-hand side above at $\mu=2$, $\tau=0$.

\begin{theorem}\label{thm-dominance}
For every even $n\ge 6$ we have $\lceil L(n)\rceil > \lceil KP(n)\rceil$.
\end{theorem}

\begin{proof}
Cross-multiplying gives the exact identities
\[
  L-KP=\frac{2^n(2n-12)}{n(n+4)(3n+2)}\ \ (n\equiv 0 \bmod 4),
  \qquad
  L-KP=\frac{2^n(2n-4)}{n(n+2)(3n+2)}\ \ (n\equiv 2\bmod 4),
\]
so the sign is decided by $2n-12$ and $2n-4$ respectively, and $L>KP$ for every
even $n\ge 6$. (At $n=4$ the number $KP$ is the larger of the two, but both
ceilings equal $7$, so no integer bound is lost.)

The inequality between real numbers does not immediately give the inequality
between ceilings, and the gap is genuinely small exactly where it matters: it
equals $8/15$ at $n=6$ and $16/39$ at $n=8$, both below $1$. Those two cases are
therefore checked directly, giving $20>19$ and $60>59$. At $n=10$ the gap is
$64/15>1$. For even $n\ge 12$ it suffices to bound the smaller of the two gaps.
Since
\[
  \frac{2n-4}{n(n+2)(3n+2)}\;\ge\;\frac{2n-12}{n(n+4)(3n+2)},
\]
the residue class $n\equiv 0\pmod 4$ is the binding one. Using $2n-12\ge n/2$
(valid for $n\ge 8$) and $n(n+4)(3n+2)\le 8n^3$, which is equivalent to
$5n^2-14n-8\ge 0$ and so holds for $n\ge 4$,
\[
  L-KP \;\ge\; \frac{2^n\,(n/2)}{8n^3} \;=\; \frac{2^n}{16n^2} \;>\; 1 .
\]
Indeed $2^{11}=2048>1936=16\cdot 11^2$, and $2^n/n^2$ is strictly increasing for
$n\ge 3$ because $\bigl(2^{n+1}/(n+1)^2\bigr)\big/\bigl(2^n/n^2\bigr)
=2n^2/(n+1)^2>1$ there. A gap exceeding $1$ forces
$\lceil L\rceil\ge L>KP+1>\lceil KP\rceil$.
\end{proof}

Both bounds are asymptotically $2^{\,n+1}/n$, so their ratio tends to $1$ and
the improvement may look negligible. It is not. The \emph{additive} gap is
asymptotically $(2/3)\cdot 2^n/n^2$, which is unbounded, so the improvement does
not run out; over the even $n$ in $[6,200]$ it survives the ceiling at all $98$
values. Table \ref{tab-1} records the tabulated positions.

\begin{table}[h]
\centering
\begin{tabular}{r|r|r|r}
$n$ & best published lower bound & Theorem \ref{thm-main} & gain \\ \hline
$6$  & $20$ (exact value, \cite{Seuranen-2007}) & $\mathbf{20}$ & $0$, matched without search \\
$8$  & $59$ \cite{KP-2021} & $\mathbf{60}$ & $+1$ \\
$10$ & $188$ \cite{Seuranen-2007,KP-2021} & $\mathbf{192}$ & $+4$ \\
$12$ & $640$ \cite{KP-2021} & $\mathbf{647}$ & $+7$ \\
$14$ & $2195$ \cite{KP-2021} & $\mathbf{2235}$ & $+40$ \\
$16$ & $7783$ \cite{KP-2021} & $\mathbf{7865}$ & $+82$
\end{tabular}
\caption{Lower bounds on $K(n,1,2)$ at the tabulated even lengths.}
\label{tab-1}
\end{table}

The entries for $n\ge 8$ are exactly the $\mu=2$ positions printed in
\cite{KP-2021}, and each printed lower bound there equals the value of their
formula, so no tabulated entry is stronger than the formula itself.

\section{Consequences at Small $n$}\label{sec-small}

\textbf{The value $K(6,1,2)=20$.} Theorem \ref{thm-main} gives
$\lceil 384/20\rceil=20$. Combined with the known twenty-word twofold covering
of $\mathbb{F}_2^6$, recorded explicitly in \cite{OEIS-A004045} as
\[
  \{0,1,2,4,8,23,27,29,30,31,39,43,45,46,47,48,49,50,52,56\}
\]
(written as integers in $[0,63]$), this recovers $K(6,1,2)=20$. The value was
previously established computationally \cite{Seuranen-2007}; Theorem
\ref{thm-main} replaces the lower-bound half of that determination by a
one-line argument, while the upper half still rests on the construction above.

\textbf{The bound $K(8,1,2)\ge 60$.} Theorem \ref{thm-main} gives
$\lceil 1536/26\rceil = 60$, one above the published $59$. The sequence
\texttt{A004045} in the On-Line Encyclopedia of Integer Sequences
\cite{OEIS-A004045} currently lists exact values only through $n=7$, namely
$2,3,4,8,12,20,32$, so $n=8$ is its first unresolved position. Krotov and
Potapov record $59\le K(8,1,2)\le 64$ \cite{KP-2021}, the upper bound being the
doubling $K(8,1,2)\le 2K(7,1,2)=64$ already present in the table of
\cite{HHKL-1993}. Theorem \ref{thm-main} raises the lower endpoint, giving
$60\le K(8,1,2)\le 64$.

\section{Even Multiplicities}\label{sec-mu}

Carry a general $\mu$ through the proof. For even $n$ the ball sum
$(n+1)+2N_2(x)$ at a codeword is odd, so parity alone forces it to exceed the
baseline $\mu(n+1)$ exactly when $\mu(n+1)$ is even --- that is, exactly when
$\mu$ is even.

\begin{proposition}\label{prop-mu}
Let $n\ge 2$ be even and let $\mu$ be an even integer with $2\le\mu\le n$.
(Since $|B(x)|=n+1$ and $C$ is a set, no $\mu$-fold covering exists for
$\mu>n+1$, so $\mu=n$ is the largest feasible even multiplicity.) Then
\[
  K(n,1,\mu)\;\ge\;\frac{\mu(\mu+1)2^n}{(\mu+1)(n+1)-1},
  \qquad\text{and hence}\qquad
  K(n,1,\mu)\;\ge\;\left\lceil\frac{\mu(\mu+1)2^n}{(\mu+1)(n+1)-1}\right\rceil .
\]
\end{proposition}

\begin{proof}
Put $g_\mu(y)=c(y)-\mu$ and $E_\mu=\sum_y g_\mu(y)=(n+1)|C|-\mu 2^n$, and let
$S_\mu=\{y:g_\mu(y)\ge 1\}$ with $s_\mu=|S_\mu|\le E_\mu$. For $n$ and $\mu$
both even, $\mu(n+1)$ is even while the ball sum $(n+1)+2N_2(x)$ at a codeword
is odd, so that sum exceeds $\mu(n+1)$ and each codeword ball meets $S_\mu$.
Counting incidences as in Step 2,
\[
  |C|\;\le\;\sum_{y\in S_\mu} c(y)\;=\;E_\mu+\mu s_\mu\;\le\;(\mu+1)E_\mu
      \;=\;(\mu+1)\bigl((n+1)|C|-\mu 2^n\bigr),
\]
which rearranges to the stated bound. \qedhere
\end{proof}

At $\mu=2$ this is Theorem \ref{thm-main}. For larger even $\mu$ the comparison
with \cite{KP-2021} depends on both parameters. Writing $P(n,\mu)$ for the
right-hand side of Proposition \ref{prop-mu} and $D=(\mu+1)(n+1)-1$, the exact
differences are
\[
  P-KP=\frac{\mu 2^n(n-3\mu)}{D\,n(n+4)}\ \ (n\equiv 0\bmod 4),
  \qquad
  P-KP=\frac{\mu 2^n(n-\mu)}{D\,n(n+2)}\ \ (n\equiv 2\bmod 4),
\]
so that, as a comparison of the underlying real-valued bounds,
\[
  P>KP \iff \mu<n/3 \quad (n\equiv 0\bmod 4),
  \qquad
  P>KP \iff \mu<n \quad (n\equiv 2\bmod 4).
\]
The gain is therefore not confined to $\mu=2$. At $\mu=4$, for instance,
Proposition \ref{prop-mu} gives $380$ at $n=10$ and $15604$ at $n=16$, against
$376$ and $15565$ from \cite{KP-2021}; it loses at $n=8$, where it gives $117$
against $118$. As in Section \ref{sec-dominance}, a strict inequality between
the real bounds need not survive the ceiling: at $n=6$, $\mu=4$ the difference
is $16/51>0$, yet both ceilings equal $38$. We therefore state the
classification for the real-valued bounds only, and do not claim a
corresponding classification after rounding.

\section{Relation to Previous Work}\label{sec-prior}

The covering excess method for \emph{ordinary} coverings goes back to van Wee
\cite{vanWee-1988}, whose radius-one bound is asymptotically $2^n/n$; that work
concerns $K(n,R)$ and does not yield the present bound. The extension of excess
counting to multiple coverings is \cite{HHKL-1993}. In their Theorem 6 the
correction term at radius one is
\[
  \varepsilon=2\left\lceil\frac{\mu(n+1)}{2}\right\rceil-\mu(n+1),
\]
so $\varepsilon=0$ exactly when $\mu(n+1)$ is even. For even $n$ the integer
$n+1$ is odd, so this happens exactly for even $\mu$, and their bound then
reduces to the sphere covering bound $K(n,1,\mu)\ge\mu 2^n/(n+1)$. The present
argument instead uses the parity of the ball sum at a codeword, which is odd
whenever $n$ is even for every $\mu$, and therefore applies precisely where
$\varepsilon$ vanishes.

Reference \cite{HHKL-1993} also cites a forthcoming manuscript of W.\ Chen and
D.\ Li entitled \emph{Lower bounds for multiple covering codes}. The author has
not located a published version; the novelty comparison above is therefore
limited to the accessible published literature.

\section*{Acknowledgment}

The author used generative language models --- Anthropic Claude (Opus 5 and
Fable 5.1), OpenAI GPT-5.6, and DeepSeek v4 --- substantially in developing and
preparing this work, for literature exploration, for generating candidate
arguments, and for drafting text. The author independently checked the
mathematical arguments, the factual claims and the bibliographic references,
made all final scholarly decisions, and takes full responsibility for the
manuscript. No financial support was received for this work, and the author
declares no conflict of interest.

%%%%%%%%%%%%%%%%%%%%%%%%%%%%%%%%%%%%%%%%%%%%%%%%%%%%%%%%%%%%%%%%%%%%
\footnotesize

\begin{thebibliography}{00}

\bibitem{HHKL-1993} H. O. H\"am\"al\"ainen, I. S. Honkala, M. K. Kaikkonen, S. Litsyn, Bounds for binary multiple covering codes, {\it Des. Codes Cryptogr.} {\bf 3} (1993) 251--275.

\bibitem{KP-2021} D. S. Krotov, V. N. Potapov, On multifold packings of radius-1 balls in Hamming graphs, {\it IEEE Trans. Inform. Theory} {\bf 67} (2021) 3585--3598.

\bibitem{Seuranen-2007} E. A. Seuranen, New lower bounds for multiple coverings, {\it Des. Codes Cryptogr.} {\bf 45} (2007) 91--94.

\bibitem{OEIS-A004045} The On-Line Encyclopedia of Integer Sequences, https://oeis.org/A004045.

\bibitem{vanWee-1988} G. J. M. van Wee, Improved sphere bounds on the covering radius of codes, {\it IEEE Trans. Inform. Theory} {\bf 34} (1988) 237--245.

\end{thebibliography}
\end{document}
